\documentclass[10pt,a4paper]{amsart}
\usepackage[margin=19mm]{geometry}
\usepackage{amsmath,amssymb,mathtools}
\usepackage[scr=rsfs,cal=euler]{mathalfa}
\usepackage{xfrac}
\usepackage[unicode,hidelinks,bookmarks=false]{hyperref}
\newcommand{\ie}{i.e.\ }
\newcommand{\cf}{cf.\ }
\newcommand{\resp}{respectively\ }
\newcommand{\Subsection}{\subsection}
\providecommand{\nobreakdash}{\nobreak\hbox{-}}
\setlength{\emergencystretch}{2em}
\multlinegap0pt
\usepackage{mdframed}
\usepackage{mdframed}
\usepackage{mdframed}
\usepackage{mdframed}
\usepackage{mdframed}
\usepackage{csquotes}
\usepackage{amssymb}
\usepackage{dsfont}
\usepackage{xcolor}
\usepackage{braket}
\usepackage[shortlabels]{enumitem}
\usepackage{mathtools}
\usepackage{cancel}
\newtheorem{thm}{}
\newcommand{\Ad}[1]{\operatorname{Ad}_{#1}}
\newcommand{\Id}{\mathds{1}}
\newcommand{\ZZ}{\mathbb{Z}}
\newcommand{\br}[1]{\left(#1\right)}
\newcommand{\calA}{\mathcal{A}}
\newcommand{\calH}{\mathcal{H}}
\newcommand{\calN}{\mathcal{N}}
\newcommand{\calU}{\mathcal{U}}
\newcommand{\ee}{\operatorname{e}}
\newcommand{\eq}[1]{\begin{align*}#1\end{align*}}
\newcommand{\eql}[1]{\begin{align}#1\end{align}}
\newcommand{\ii}{\operatorname{i}}
\newcommand{\ip}[2]{\langle #1, #2 \rangle}
\newcommand\norm[1]{\left\lVert#1\right\rVert}
\newcommand{\ve}{\varepsilon}
\title{The $\ZZ_2$-Index of a Pair of Pure States and the\\Topology of Interacting 1D Superconductors}
\author{Anna Mazhar, Jacob Shapiro}
\date{Source: arXiv:2608.28796v1 (CC BY 4.0)}
\begin{document}
\maketitle

\noindent\emph{Source and licence.} The $\ZZ_2$-Index of a Pair of Pure States and the\\Topology of Interacting 1D Superconductors. Version arXiv:2608.28796v1 (CC BY 4.0), distributed by arXiv under the Creative Commons Attribution 4.0 International licence, http://creativecommons.org/licenses/by/4.0/, as recorded in the arXiv licence metadata for this exact version and shown on the abstract page https://arxiv.org/abs/2608.28796v1. Full licence notice: \texttt{licenses/arxiv-2608.28796-CC-BY-4.0.txt}.\par\smallskip

\noindent\textit{Editorial scope.} Counted unit: the article's thm thm (source0.tex lines 7518--7522). The article's complete original proof of it is delivered with it (source0.tex lines 7524--7955, one \texttt{proof} environment of 432 lines). No mathematics is written, repaired or re-derived: the statement and the proof are the article's own lines, in the article's own notation, and no retained line is substituted or re-flowed.  No further source-local context is needed: every cross-reference of every retained line resolves inside the two delivered spans.  Nothing was omitted from any retained span, so the package carries no omission record.  No retained line contains a citation, so the wrapper carries no bibliography and no bibliography entry.  No retained line contains a figure environment, an image-inclusion command or drawing code, so no image is delivered and the package declares no asset dependency.  Editorial formatting only, in addition: an AMS \texttt{amsart} preamble that restates the article's own declarations and macros used by the retained lines, namely the environments \texttt{thm} on separate counters, together with the standard packages those lines need.  The retained environments print in this document as Theorem 1 in order of appearance, whereas the article prints the same items with the numbering of its own full document.  The complete native source of the article as distributed in the arXiv e-print for this exact version is bundled unchanged as \texttt{sources/208811/source0.tex}.
\begin{thm}[LGA-injectivity of the index] Let $\omega=\omega_{\rm p}\circ\beta_1$ be a $\theta$-invariant SRE state with a $\theta$-symmetric generator for its LGA $t\mapsto\beta_t$, such that $\calN(\omega_{\rm p},\omega_0)=+1$ where $\omega_0$ is the $k=0$ Kitaev state.

Then there exists an LGA $t\mapsto \alpha_t$ with a $\theta$-symmetric generator such that 
\eql{\omega = \omega_0\circ\alpha_1\,.} 
\end{thm}

\begin{proof}
Since $P_0$ has no matrix elements coupling distinct sites, the $k=0$
Kitaev state $\omega_0$ is a pure product state.

For $j\in\Set{{\rm p},0}$, write
\eq{
\omega_j
=
\hat{\bigotimes}_{x\in\ZZ}
\omega_{j,x}.
}
Let $\Gamma_x\in\calA_{\Set{x}}$ denote the one-site parity unitary, so
that
\eq{
\Gamma_x^\ast
&=
\Gamma_x,\\
\Gamma_x^2
&=
\Id,\\
\theta(a)
&=
\Gamma_xa\Gamma_x
\qquad
\br{
a\in\calA_{\Set{x}}
}.
}
Since $\omega_{j,x}$ is pure and $\theta$-invariant, it is represented by
a unit vector $\xi_{j,x}$ of definite local parity. Thus there is some
$\ve_{j,x}\in\Set{\pm1}$ such that
\eq{
\Gamma_x\xi_{j,x}
=
\ve_{j,x}\xi_{j,x}.
}
Set
\eq{
D
:=
\Set{
x\in\ZZ:
\ve_{{\rm p},x}\neq\ve_{0,x}
}.
}

We first show that $D$ is finite. Since
$\calN(\omega_{\rm p},\omega_0)$ is defined, there is a unitary
$u\in\calU(\calA)$ such that
\eq{
\omega_0
=
\omega_{\rm p}\circ\Ad{u}.
}
Choose a finite set $\Lambda\subseteq\ZZ$ and a local unitary
$v\in\calU(\calA_\Lambda)$ satisfying
\eq{
\norm{u-v}
<
1.
}
If $x\notin\Lambda$, then $\Gamma_x$ is even and has support disjoint
from that of $v$, so
\eq{
v^\ast\Gamma_xv
=
\Gamma_x.
}
Consequently,
\eq{
\left|
\ve_{0,x}-\ve_{{\rm p},x}
\right|
&=
\left|
\omega_{\rm p}
\br{
u^\ast\Gamma_xu
}
-
\omega_{\rm p}
\br{
v^\ast\Gamma_xv
}
\right|\\
&\leq
\norm{
u^\ast\Gamma_xu
-
v^\ast\Gamma_xv
}\\
&\leq
2\norm{u-v}\\
&<
2.
}
Since $\ve_{0,x}$ and $\ve_{{\rm p},x}$ both belong to $\Set{\pm1}$, this
implies
\eq{
\ve_{0,x}
=
\ve_{{\rm p},x}
\qquad
\br{
x\notin\Lambda
}.
}
Thus
\eq{
D
\subseteq
\Lambda,
}
and in particular $D$ is finite.

We next relate the parity of $\lvert D\rvert$ to the relative index.
Let
\eq{
\br{
\pi_{\rm p},
\calH_{\rm p},
\Omega_{\rm p}
}
}
be the GNS representation of $\omega_{\rm p}$, and let
$\Theta_{\rm p}$ be the unitary implementing $\theta$, normalized by
\eq{
\Theta_{\rm p}\Omega_{\rm p}
=
\Omega_{\rm p}.
}
In the product realization of this representation,
\eq{
\Omega_{\rm p}
=
\hat{\bigotimes}_{x\in\ZZ}
\xi_{{\rm p},x},
}
and
\eq{
\Theta_{\rm p}
=
\hat{\bigotimes}_{x\in\ZZ}
\br{
\ve_{{\rm p},x}\Gamma_x
}.
}
Indeed, each local factor $\ve_{{\rm p},x}\Gamma_x$ fixes
$\xi_{{\rm p},x}$ and implements the local parity automorphism.

Set
\eq{
\Omega_0
:=
\pi_{\rm p}(u)\Omega_{\rm p}.
}
Then $\Omega_0$ represents $\omega_0$ in the same GNS representation.
Since $\omega_0$ is a pure product state, the standard incomplete
tensor-product description gives, up to an overall phase,
\eq{
\Omega_0
=
\hat{\bigotimes}_{x\in\ZZ}
\xi_{0,x}.
}
It follows that
\eq{
\Theta_{\rm p}\Omega_0
&=
\hat{\bigotimes}_{x\in\ZZ}
\br{
\ve_{{\rm p},x}\Gamma_x\xi_{0,x}
}\\
&=
\br{
\prod_{x\in\ZZ}
\ve_{{\rm p},x}\ve_{0,x}
}
\Omega_0\\
&=
\br{
-1
}^{\lvert D\rvert}
\Omega_0.
}
On the other hand,
\eq{
\calN(\omega_{\rm p},\omega_0)
&=
\omega_{\rm p}
\br{
u^\ast\theta(u)
}\\
&=
\ip{
\pi_{\rm p}(u)\Omega_{\rm p}
}{
\Theta_{\rm p}\pi_{\rm p}(u)\Omega_{\rm p}
}\\
&=
\ip{
\Omega_0
}{
\Theta_{\rm p}\Omega_0
}.
}
Therefore,
\eql{
\calN(\omega_{\rm p},\omega_0)
=
\br{
-1
}^{\lvert D\rvert}.
}
The hypothesis
\eq{
\calN(\omega_{\rm p},\omega_0)
=
1
}
thus implies that $\lvert D\rvert$ is even.

We now construct a symmetric LGA connecting $\omega_0$ to
$\omega_{\rm p}$. For every $x\notin D$, the vectors $\xi_{0,x}$ and
$\xi_{{\rm p},x}$ belong to the same parity eigenspace. Hence there is
an even one-site unitary
\eq{
U_x
\in
\calU\br{
\calA_{\Set{x}}
}
}
such that
\eq{
U_x\xi_{0,x}
&=
\xi_{{\rm p},x},\\
\theta(U_x)
&=
U_x.
}
Since the even unitary group of the finite-dimensional algebra
$\calA_{\Set{x}}$ is connected, we may write
\eq{
U_x
=
\ee^{-\ii h_x},
}
where
\eq{
h_x^\ast
&=
h_x,\\
\theta(h_x)
&=
h_x,\\
\norm{h_x}
&\leq
\pi.
}
Define the time-independent zero-chain
\eq{
\Phi^{(1)}_t(x)
:=
\begin{cases}
h_x,
&
x\notin D,\\
0,
&
x\in D.
\end{cases}\,.
}
This zero-chain is uniformly bounded, strictly on-site, and
$\theta$-even. Let $t\mapsto\tau_t$ be the corresponding LGA. Since the
terms at distinct sites are even and have disjoint supports, they
commute, and hence
\eq{
\omega_0\circ\tau_1
=
\br{
\hat{\bigotimes}_{x\notin D}
\omega_{{\rm p},x}
}
\hat{\otimes}
\br{
\hat{\bigotimes}_{x\in D}
\omega_{0,x}
}.
}

It remains to change the finitely many factors indexed by $D$. Since
$\lvert D\rvert$ is even, choose a pairing
\eq{
D
=
\Set{
x_1,y_1,\ldots,x_m,y_m
}.
}
For each $r\in\Set{1,\ldots,m}$, we have
\eq{
\ve_{{\rm p},x_r}
&=
-\ve_{0,x_r},\\
\ve_{{\rm p},y_r}
&=
-\ve_{0,y_r},
}
and therefore
\eq{
\ve_{{\rm p},x_r}\ve_{{\rm p},y_r}
=
\ve_{0,x_r}\ve_{0,y_r}.
}
Thus the vectors
\eq{
\xi_{0,x_r}
\hat{\otimes}
\xi_{0,y_r}
}
and
\eq{
\xi_{{\rm p},x_r}
\hat{\otimes}
\xi_{{\rm p},y_r}
}
belong to the same total-parity eigenspace of the two-site algebra
$\calA_{\Set{x_r,y_r}}$. Consequently, there is an even unitary
\eq{
V_r
\in
\calU\br{
\calA_{\Set{x_r,y_r}}
}
}
such that
\eq{
V_r
\br{
\xi_{0,x_r}
\hat{\otimes}
\xi_{0,y_r}
}
=
\xi_{{\rm p},x_r}
\hat{\otimes}
\xi_{{\rm p},y_r}.
}
As before, we may write
\eq{
V_r
=
\ee^{-\ii k_r},
}
where
\eq{
k_r^\ast
&=
k_r,\\
\theta(k_r)
&=
k_r,\\
\norm{k_r}
&\leq
\pi.
}

Define a second time-independent zero-chain by
\eq{
\Phi^{(2)}_t(x)
:=
\begin{cases}
k_r,
&
x=x_r
\text{ for some }
r\in\Set{1,\ldots,m},\\
0,
&
\text{otherwise}.
\end{cases}\,.
}
Since only finitely many terms occur, this is a finite-range,
$\theta$-even zero-chain. Let $t\mapsto\rho_t$ be its LGA. The pairs are
disjoint, and all the terms are even, so
\eq{
\br{
\omega_0\circ\tau_1
}
\circ\rho_1
=
\omega_{\rm p}.
}

After a continuous time reparametrization, we may concatenate
$t\mapsto\tau_t$ and $t\mapsto\rho_t$ to obtain an LGA
$t\mapsto\gamma_t$ with a $\theta$-even generator and endpoint
\eq{
\gamma_1
=
\tau_1\circ\rho_1.
}
It satisfies
\eql{
\omega_{\rm p}
=
\omega_0\circ\gamma_1.
}

Finally, concatenate $t\mapsto\gamma_t$ with the given LGA
$t\mapsto\beta_t$. Since both generators are $\theta$-even, the
concatenated evolution $t\mapsto\alpha_t$ also has a $\theta$-even
generator, and its endpoint is
\eq{
\alpha_1
=
\gamma_1\circ\beta_1.
}
Therefore,
\eq{
\omega_0\circ\alpha_1
&=
\omega_0\circ\gamma_1\circ\beta_1\\
&=
\omega_{\rm p}\circ\beta_1\\
&=
\omega.
}
This proves the claim.
\end{proof}

\end{document}
